\chapter{Lists, virtualisation, and infinite scroll}
\label{ch:lists}

\begin{aside}
A list with a million items must not paint a million rows.
Virtualisation is the trick: render only what's visible, fake
the rest with scroll offsets.
\end{aside}

\section{The viewport}

A list has a viewport (visible rows) and a content (total
rows). When content is larger than viewport, scrolling moves
a sliding window.

\begin{lstlisting}
class VirtualList {
    int viewport_height;
    int total_rows;
    int scroll_top;
public:
    int first_visible() { return scroll_top; }
    int last_visible() {
        return scroll_top + viewport_height;
    }
};
\end{lstlisting}

Only rows in $[\text{first}, \text{last})$ are rendered.

\section{Row height}

Fixed-height rows make the math trivial: row $i$ is at $i
\times h$. Variable heights need a prefix sum or a
binary-searchable accumulator.

\maya{}'s list defaults to fixed height; variable-height
needs a measure callback per row.

\section{Item factory}

The list takes a factory that returns the element for an
index:

\begin{lstlisting}
list(total_rows, [](int i) {
    return text("row " + std::to_string(i));
});
\end{lstlisting}

The list calls the factory only for visible rows.

\section{Scroll position}

A scroll signal drives the viewport:

\begin{lstlisting}
auto scroll = signal(0);
auto list = virtual_list(items, scroll);
\end{lstlisting}

Keyboard handlers update the signal: \code{j} increments,
\code{k} decrements, page-down jumps by viewport.

\section{Smooth scrolling}

Snapping to the next row feels abrupt. Smooth scrolling
animates the scroll signal:

\begin{lstlisting}
animate_to(scroll, target, 150ms);
\end{lstlisting}

The viewport slides over a tween instead of jumping.

\section{Infinite scroll}

When the user scrolls near the bottom, fetch more items:

\begin{lstlisting}
ctx.effect([&] {
    if (ctx.get(scroll) > total_rows - viewport - 10) {
        load_more();
    }
});
\end{lstlisting}

\code{load\_more} appends to the items signal; the list
extends.

\section{Sticky headers}

A scrolling list with section headers can pin the current
section's header to the top:

\begin{lstlisting}
auto pinned_header = memo([&] {
    return section_at(ctx.get(scroll));
});
\end{lstlisting}

Render the pinned header above the viewport.

\section{Lazy loading}

Items don't need to be in memory --- a windowed data
source fetches as the viewport approaches. \maya{} provides
\code{LazyList} that calls a fetch callback for unloaded
ranges.

\section{Recap}

\begin{recap}
\begin{itemize}[leftmargin=*]
  \item Render only visible rows.
  \item Fixed height: trivial math; variable: prefix sums.
  \item Scroll signal drives the viewport.
  \item Smooth scroll animates; infinite scroll fetches on
    edge.
  \item Sticky headers via memo on scroll position.
\end{itemize}
\end{recap}

\section*{Workshop}

\begin{enumerate}[leftmargin=*]
  \item Build a list of 100,000 rows. Confirm only ~50 are
    rendered at a time.
  \item Add smooth scroll with a 150 ms tween.
  \item Implement infinite scroll for a paginated API.
\end{enumerate}
